% ---------------------------------------------------------------
%  Curriculum vitae — Abdul Razak
%  Engine: tectonic (XeTeX).  Build:  tectonic -X compile cv.tex
% ---------------------------------------------------------------
\documentclass[9pt,letterpaper]{extarticle}

\usepackage[T1]{fontenc}
\usepackage[utf8]{inputenc}
\usepackage{newpxtext}                 % Palatino-style text face
\usepackage{microtype}
\usepackage[
  top=0.62in, bottom=0.58in, left=0.74in, right=0.74in,
  footskip=0.26in
]{geometry}
\usepackage{enumitem}
\usepackage{titlesec}
\usepackage{xcolor}
\usepackage{fancyhdr}
\usepackage[hidelinks]{hyperref}

\definecolor{ink}{HTML}{111820}
\definecolor{accent}{HTML}{143352}
\definecolor{soft}{HTML}{6B7684}
\definecolor{rule}{HTML}{C9D0DA}

\hypersetup{
  colorlinks=true,
  urlcolor=accent,
  linkcolor=accent,
  pdftitle={Abdul Razak --- Curriculum Vitae},
  pdfauthor={Abdul Razak},
  pdfsubject={Curriculum Vitae},
  pdfkeywords={artificial intelligence, multi-objective decision making,
               probabilistic graphical models, generative AI}
}

\color{ink}
\setlength{\parindent}{0pt}
\linespread{1.03}

% ---- section headings ------------------------------------------------
\titleformat{\section}
  {\normalfont\footnotesize\bfseries\scshape\color{accent}}
  {}{0pt}{\MakeUppercase}[{\color{rule}\titlerule[0.5pt]}]
\titlespacing*{\section}{0pt}{7.5pt plus 2pt minus 2pt}{3.5pt}

% ---- date-gutter list ------------------------------------------------
\newlength{\gutter}\setlength{\gutter}{0.92in}
\newenvironment{dated}[1][5pt]
  {\begin{list}{}{%
      \setlength{\labelwidth}{\gutter}
      \setlength{\labelsep}{0.18in}
      \setlength{\leftmargin}{\dimexpr\gutter+0.18in\relax}
      \setlength{\itemsep}{#1}
      \setlength{\parsep}{0pt}
      \setlength{\topsep}{3pt}
      \renewcommand{\makelabel}[1]{\raggedright\footnotesize\color{soft}##1\hfill}}}
  {\end{list}}

\newcommand{\role}[2]{{\normalsize\bfseries #1}\\[1pt]{\color{accent}#2}}
\newcommand{\degree}[2]{{\normalsize\bfseries #1}\\[1pt]{\color{accent}#2}}

\newenvironment{duties}
  {\begin{itemize}[leftmargin=11pt, topsep=2pt, itemsep=1.2pt,
                   parsep=0pt, label={\color{rule}\rule[2.2pt]{4pt}{0.5pt}}]}
  {\end{itemize}}

% ---- reference lists (publications, patents) -------------------------
\newenvironment{refs}
  {\begin{list}{}{%
      \setlength{\labelwidth}{0.34in}
      \setlength{\labelsep}{0.10in}
      \setlength{\leftmargin}{0.44in}
      \setlength{\itemsep}{2.8pt}
      \setlength{\parsep}{0pt}
      \setlength{\topsep}{3pt}
      \renewcommand{\makelabel}[1]{\footnotesize\color{accent}##1\hfill}}}
  {\end{list}}

\newcommand{\me}{\textbf{A.\,Razak}}
\newcommand{\venue}[1]{\textit{\color{soft}#1}}
\newcommand{\group}[1]{\smallskip{\footnotesize\scshape\color{soft}#1}\par\smallskip}

\newenvironment{plainlist}
  {\begin{itemize}[leftmargin=11pt, topsep=2pt, itemsep=2.2pt, parsep=0pt,
                   label={\color{accent}\rule[1.6pt]{2.6pt}{2.6pt}}]}
  {\end{itemize}}

% ---- running footer --------------------------------------------------
\pagestyle{fancy}
\fancyhf{}
\renewcommand{\headrulewidth}{0pt}
\renewcommand{\footrulewidth}{0.4pt}
\renewcommand{\footrule}{\hbox to\headwidth{\color{rule}\leaders\hrule height \footrulewidth\hfill}}
\fancyfoot[L]{\footnotesize\color{soft}Abdul Razak}
\fancyfoot[R]{\footnotesize\color{soft}%
  Detailed project notes: \href{https://abdulrazakucc.github.io/}{online portfolio}%
  \quad\textbar\quad Page \thepage\ of 2}

\begin{document}

% ======================= masthead =====================================
\begin{center}
  {\fontsize{22}{24}\selectfont\scshape Abdul Razak}\\[4pt]
  {\normalsize\itshape Staff Data Scientist, Extreme Networks}\\[7pt]
  {\color{rule}\rule{\textwidth}{0.6pt}}\\[5pt]
  {\footnotesize
   Salem, United States \quad\textbar\quad
   \href{https://abdulrazakucc.github.io/}{Portfolio} \quad\textbar\quad
   \href{https://scholar.google.com/citations?user=NeuA3uAAAAAJ\&hl=en}{Google Scholar} \quad\textbar\quad
   \href{https://dblp.org/pid/120/1333.html}{DBLP} \quad\textbar\quad
   \href{https://orcid.org/0009-0004-6819-147X}{ORCID} \quad\textbar\quad
   \href{https://patents.justia.com/inventor/abdul-razak}{Patents} \quad\textbar\quad
   \href{https://www.linkedin.com/in/data-analytics/}{LinkedIn}}
\end{center}
\vspace{2pt}

% ======================= research =====================================
\section{Research}

My work concerns decision making when objectives conflict and preferences are known only in part.
During my doctorate and post-doctoral fellowship I developed exact and approximate algorithms for
multi-objective constraint optimisation and influence diagrams, and methods that elicit only as much
preference information as a decision actually requires; this appeared at AAAI, IJCAI, UAI, CP and
ICTAI, largely with colleagues at IBM Research Dublin. In industry I have applied the same
probabilistic and decision-theoretic tools to building security, where the work produced eight
granted patents, and since 2020 to enterprise wireless networks, where I lead the data-science work
on generative models for network diagnosis.

\smallskip
{\footnotesize\textsc{\color{soft}Interests}}\; Multi-objective decision making under uncertainty;
preference elicitation and inference; probabilistic graphical models and influence diagrams;
grounded and retrieval-augmented generation; constraint optimisation; machine learning for
networked systems.

% ======================= education ====================================
\section{Education}
\begin{dated}
  \item[2010--2014]
    \degree{PhD, Computer Science (Artificial Intelligence)}
           {National University of Ireland, Cork --- Cork Constraint Computation Centre}\\[1.5pt]
    Thesis: \textit{Reasoning with Imprecise Trade-offs in Decision Making under Certainty and
    Uncertainty}. Supervised by Nic Wilson. Funded by the IRCSET\,/\,IBM Enterprise Partnership
    Scheme as part of IBM's Risk Management Colloboratory.
  \item[2003--2005]
    \degree{MSc, Mathematics with Computer Science}{Osmania University, Hyderabad, India}\\[1.5pt]
    Awarded the university merit scholarship.
  \item[2000--2003]
    \degree{BSc, Mathematics, Statistics and Computer Science}{Osmania University, Hyderabad, India}
\end{dated}

% ======================= appointments =================================
\section{Appointments}
\begin{dated}
  \item[2020--present]
    \role{Staff Data Scientist}{Extreme Networks --- Salem, United States; previously Cork, Ireland}
    \begin{duties}
      \item Lead the data-science work on generative models for network diagnosis. Language models
            are used for hypothesis generation and explanation; anomaly detection, a knowledge graph
            of the deployment and live telemetry supply the evidence they reason over.
      \item Built agent workflows that triage client connectivity, roaming and device-health
            problems across deployments of thousands of access points and switches, and draft
            remediation steps for network operators.
      \item Developed correlation and inference pipelines over telemetry and client-experience
            metrics, so that investigation of routine faults takes seconds rather than hours.
      \item Adapted the preference-inference methods from my research to rank recommendations for
            customers whose priorities differ in reliability, cost and coverage.
    \end{duties}
  \item[2017--2020]
    \role{Data Scientist}{Johnson Controls --- Cork, Ireland}
    \begin{duties}
      \item Machine-learning lead for the intrusion-alarm, access-control and retail-analytics
            product lines.
      \item Alarm-panel classification by Pareto optimisation, root-cause analysis of events, and
            site-independent signature generation for predictive security analysis; eight granted
            US patents issued from this work.
      \item Relational probabilistic models in BLOG and Bayesian inference in PyMC3, deployed
            across large multi-site estates.
    \end{duties}
  \item[2014--2017]
    \role{Post-Doctoral Fellow (Science Foundation Ireland)}
         {Insight Centre for Data Analytics, NUI Cork}
    \begin{duties}
      \item Probabilistic graphical models, clustering and influence diagrams for industrial
            decision problems, within a programme run with United Technology Research Centre.
      \item Preference inference with lexicographic and weighted-average models, partly with IBM
            Research Dublin (AAAI 2017, IJCAI 2015, ICTAI 2015).
    \end{duties}
  \item[2005--2010]
    \role{Lecturer in Mathematics}{Osmania University --- Hyderabad, India}
    \begin{duties}
      \item Undergraduate teaching: abstract and linear algebra, real and numerical analysis,
            differential and vector calculus, probability and statistics.
    \end{duties}
\end{dated}

\clearpage

% ======================= publications =================================
\section{Publications}
{\footnotesize\color{soft}Peer-reviewed; complete list on
 \href{https://scholar.google.com/citations?user=NeuA3uAAAAAJ\&hl=en}{Google Scholar} and
 \href{https://dblp.org/pid/120/1333.html}{DBLP}.}
\begin{refs}
  \item[{[C6]}] R.~Marinescu, \me, N.~Wilson. Multi-objective influence diagrams with possibly
        optimal policies. \venue{Proc. 31st AAAI Conference on Artificial Intelligence, 2017.}
  \item[{[C5]}] Anne-Marie, \me, N.~Wilson. The comparison of multi-objective preference inference
        based on lexicographic and weighted average models. \venue{Proc. 27th IEEE International
        Conference on Tools with Artificial Intelligence (ICTAI), 2015.}
  \item[{[C4]}] N.~Wilson, \me, R.~Marinescu. Computing possibly optimal solutions for
        multi-objective constraint optimisation with trade-offs. \venue{Proc. 24th International
        Joint Conference on Artificial Intelligence (IJCAI), 2015.}
  \item[{[W1]}] Anne-Marie, \me, N.~Wilson. Multi-objective constraint optimisation with
        lexicographic preference models. \venue{Workshop on Algorithmic Issues for Inference in
        Graphical Models (AIGM), 2015.}
  \item[{[C3]}] R.~Marinescu, \me, N.~Wilson. Multi-objective constraint optimisation with
        trade-offs. \venue{Proc. 19th International Conference on Principles and Practice of
        Constraint Programming (CP), 2013.}
  \item[{[C2]}] R.~Marinescu, \me, N.~Wilson. Multi-objective influence diagrams with trade-offs.
        \venue{Proc. 28th Conference on Uncertainty in Artificial Intelligence (UAI), 2012.}
\end{refs}

% ======================= patents ======================================
\section{Granted US Patents}
{\footnotesize\color{soft}Nine granted between 2019 and 2021; filings searchable on
 \href{https://patents.justia.com/inventor/abdul-razak}{Justia} and
 \href{https://patents.google.com/?inventor=Abdul+Razak}{Google Patents}.}

\group{Decision intelligence and maintenance}
\begin{refs}
  \item[{[P9]}] \me, L.\,A.~Deleris, R.~Marinescu, N.~Wilson. Interactive method to reduce the
        amount of trade-off information required from decision-maker in multi-attribute
        decision-making under uncertainty. \venue{Granted 2019.}
  \item[{[P8]}] \me, G.~Subramanian, C.\,J.~Donovan. Virtual maintenance manager.
        \venue{Granted 2019.}
  \item[{[P7]}] D.~Horgan, \me, G.~Bernardi. Root-cause analysis of event occurrences.
        \venue{Granted 2021.}
\end{refs}

\group{Alarm management}
\begin{refs}
  \item[{[P6]}] \me, M.~Stewart, G.~Subramanian. Automated alarm panel classification using Pareto
        optimization. \venue{Granted 2020.}
  \item[{[P5]}] \me. Systems and methods for managing alarm data of multiple locations.
        \venue{Granted 2021.}
  \item[{[P4]}] M.~Costello, F.~Fala, J.\,R.~Holliday, N.\,B.~Pavel, \me. Building security system
        with alarm lists and alarm rule editing. \venue{Granted 2021.}
\end{refs}

\group{Building security analytics}
\begin{refs}
  \item[{[P3]}] M.~Dziduch, C.\,J.~Donovan, \me. Building security analysis system with
        site-independent signature generation for predictive security analysis. \venue{Granted 2020.}
  \item[{[P2]}] M.~Dziduch, C.\,J.~Donovan, \me. Rules-based method of identifying misuse of
        emergency fire exits using data generated by a security alarm system. \venue{Granted 2020.}
  \item[{[P1]}] M.~Dziduch, C.\,J.~Donovan, \me. Systems and methods for managing a security system.
        \venue{Granted 2020.}
\end{refs}

% ======================= honours ======================================
\section{Honours and Funding}
\begin{plainlist}
  \item IEEE Senior Member, elevated May 2026.
  \item Science Foundation Ireland post-doctoral fellowship, Insight Centre for Data Analytics,
        NUI Cork, 2014--2017.
  \item IRCSET\,/\,IBM Enterprise Partnership Scheme scholarship; full doctoral funding, 2010--2013.
  \item Osmania University merit scholarship (MSc, 2003--2005); ranked 40th of roughly 30{,}000
        candidates in the university mathematics entrance examination.
\end{plainlist}

% ======================= service ======================================
\section{Service and Talks}
\begin{plainlist}
  \item Reviewer, \textit{Computers and Electrical Engineering} (Elsevier); more than ten
        manuscripts in artificial intelligence, wireless networking and intelligent systems.
  \item Reviewer, IEEE International Conference on Communications (ICC); artificial-intelligence
        submissions to the 2025 technical programme.
  \item Invited talks: Sapienza University of Rome, on preference handling in decision making under
        certainty and uncertainty; NUI Galway, on probabilistic decision making under uncertainty;
        NUI Cork and IBM Research Dublin, during the Risk Management Colloboratory and United
        Technology Research Centre programmes.
  \item Conference presentations, oral and poster, at AAAI, IJCAI, UAI, CP and ICTAI.
  \item Five years of undergraduate lecturing; mentoring of junior data scientists in industry.
\end{plainlist}

% ======================= skills =======================================
\section{Tools}
\begin{description}[leftmargin=0.76in, style=sameline, font={\footnotesize\scshape\color{soft}},
                    topsep=2pt, itemsep=1.6pt, parsep=0pt]
  \item[Programming] Python, C++ (STL, Boost, Eigen), SQL, PySpark, PyFlink.
  \item[Modelling] PyMC3, BLOG, Swift, ILOG CPLEX, LPSOLVE, scikit-learn.
  \item[Platforms] AWS SageMaker and Athena, Azure Databricks, streaming and batch pipelines.
  \item[Languages] English, Hindi, Telugu, Urdu. Irish citizen since 2021.
\end{description}

\end{document}
